\section{Methodology}


I manually downloaded the PDF on each canton's website.
{\footnotesize
\setlength{\tabcolsep}{3pt}
\renewcommand{\arraystretch}{1.1}

\begin{minipage}[t]{0.49\textwidth}
\centering

% === First half of table ===
\begin{longtable}{|p{1cm}|p{1cm}|p{1.5cm}|p{1.5cm}|p{1cm}|}
\hline
\rotatebox{70}{\textbf{Canton}} &
\rotatebox{70}{\textbf{Version}} &
\rotatebox{70}{\textbf{Year}} &
\rotatebox{70}{\textbf{Availability}} &
\rotatebox{70}{\textbf{Language}} \\ 
\hline

\multirow{3}{*}{AG} & 3 & 2025 & \href{https://www.ag.ch/de/themen/planen-bauen/raumentwicklung/grundlagen-und-kantonalplanung/richtplan}{Online} & DE \\
 & 2 &  & † & DE \\
 & 1 &  & ‡ & DE \\ \hline

\multirow{3}{*}{AR} & 3 & 2025 & \href{https://ar.ch/verwaltung/departement-bau-und-volkswirtschaft/amt-fuer-raum-und-wald/abteilung-raumentwicklung/planen/kantonale-planung/}{Online} & DE \\
 & 2 &  & † & DE \\
 & 1 &  & ‡ & DE \\ \hline

\multirow{3}{*}{AI} & 3 & 2018 & \href{https://www.ai.ch/themen/planen-und-bauen/raumplanung/richtplanung}{Online} & DE \\
 & 2 &  & † & DE \\
 & 1 &  & ‡ & DE \\ \hline

\multirow{3}{*}{BE} & 3 & 2024 & \href{https://www.raumplanung.dij.be.ch/fr/start/kantonaler-richtplan.html}{Online} & DE \\
 & 2 &  & † & DE \\
 & 1 &  & ‡ & DE \\ \hline

\multirow{3}{*}{BL} & 3 & 2009 & \href{https://www.baselland.ch/politik-und-behorden/direktionen/bau-und-umweltschutzdirektion/raumplanung/raumplanung/richtplanung}{Online} & DE \\
 & 2 &  & † & DE \\
 & 1 &  & ‡ & DE \\ \hline

\multirow{3}{*}{BS} & 3 & 2025 & \href{https://www.bs.ch/bvd/staedtebau-architektur/raumplanung/kantonaler-richtplan}{Online} & DE \\
 & 2 &  & † & DE \\
 & 1 &  & ‡ & DE \\ \hline

\multirow{3}{*}{FR} & 3 & 2020 & \href{https://www.fr.ch/dime/seca/plan-directeur-cantonal}{Online} & FR \\
 & 2 &  & † & FR \\
 & 1 &  & ‡ & FR \\ \hline

\multirow{3}{*}{GE} & 3 & 2024 & \href{https://www.ge.ch/dossier/amenager-territoire/planification-cantonale-regionale/plan-directeur-cantonal-2030}{Online} & FR \\
 & 2 &  & † & FR \\
 & 1 &  & ‡ & FR \\ \hline

\multirow{3}{*}{GL} & 3 & 2025 & \href{https://www.gl.ch/verwaltung/bau-und-umwelt/hochbau/raumentwicklung-und-geoinformation/raumentwicklung/kantonale-richtplanung.html/5046}{Online} & DE \\
 & 2 &  & † & DE \\
 & 1 &  & ‡ & DE \\ \hline

\multirow{3}{*}{GR} & 3 & 2025 & \href{https://www.gr.ch/DE/institutionen/verwaltung/dvs/are/dienstleistungen/richtplanung/Seiten/Text-und-Karte.aspx}{Online} & DE \\
 & 2 &  & † & DE \\
 & 1 &  & ‡ & DE \\ \hline

\multirow{3}{*}{JU} & 3 &  & \href{https://www.jura.ch/fr/Autorites/Administration/DEN/SDT/Plan-directeur-cantonal/Plan-directeur-cantonal.html}{Online} & FR \\
 & 2 &  & † & FR \\
 & 1 &  & ‡ & FR \\ \hline

\multirow{3}{*}{LU} & 3 & 2009 & \href{https://richtplan.lu.ch/}{Online} & DE \\
 & 2 &  & † & DE \\
 & 1 &  & ‡ & DE \\ \hline

\multirow{3}{*}{NE} & 3 & 2011 & \href{https://www.ne.ch/autorites/DDTE/SCAT/pdc/Pages/accueil.aspx}{Online} & FR \\
 & 2 &  & † & FR \\
 & 1 &  & ‡ & FR \\ \hline

\caption{Cantonal master plans (part 1).}
\end{longtable}

\end{minipage}
\hfill
\begin{minipage}[t]{0.49\textwidth}
\centering


% === Second half of table ===
\begin{longtable}{|p{1cm}|p{1cm}|p{1.5cm}|p{1.5cm}|p{1cm}|}
\hline
\rotatebox{70}{\textbf{Canton}} &
\rotatebox{70}{\textbf{Version}} &
\rotatebox{70}{\textbf{Year}} &
\rotatebox{70}{\textbf{Availability}} &
\rotatebox{70}{\textbf{Language}} \\ 
\hline

\multirow{3}{*}{NW} & 3 & 2002 & \href{https://www.nw.ch/raumentwdienste/2871}{Online} & DE \\
 & 2 &  & † & DE \\
 & 1 &  & ‡ & DE \\ \hline

\multirow{3}{*}{OW} & 3 & 2019 & \href{https://www.ow.ch/dienstleistungen/1527}{Online} & DE \\
 & 2 &  & † & DE \\
 & 1 &  & ‡ & DE \\ \hline

\multirow{3}{*}{SG} & 3 & 2017 & \href{https://www.sg.ch/bauen/raumentwicklung/kantonaleplanung/richtplanung.html}{Online} & DE \\
 & 2 &  & † & DE \\
 & 1 &  & ‡ & DE \\ \hline

\multirow{3}{*}{SH} & 3 & 2014 & \href{https://sh.ch/CMS/Webseite/Kanton-Schaffhausen/Beh-rde/Verwaltung/Baudepartement/Planungs--und-Naturschutzamt/Raumplanung/Kantonale-Richtplanung-1612737-DE.html}{Online} & DE \\
 & 2 &  & † & DE \\
 & 1 &  & ‡ & DE \\ \hline

\multirow{3}{*}{SZ} & 3 & 2003 & \href{https://www.sz.ch/behoerden/verwaltung/volkswirtschaftsdepartement/amt-fuer-raumentwicklung/kantonaler-richtplan.html/8756-8758-8802-10373-11039-12578}{Online} & DE \\
 & 2 &  & † & DE \\
 & 1 &  & ‡ & DE \\ \hline

\multirow{3}{*}{SO} & 3 & 2017 & \href{https://so.ch/verwaltung/bau-und-justizdepartement/amt-fuer-raumplanung/richtplanung/}{Online} & DE \\
 & 2 &  & † & DE \\
 & 1 &  & ‡ & DE \\ \hline

\multirow{3}{*}{TI} & 3 & 2013 & \href{https://www4.ti.ch/dt/dstm/sst/temi/piano-direttore/piano-direttore/}{Online} & IT \\
 & 2 &  & † & IT \\
 & 1 &  & ‡ & IT \\ \hline

\multirow{3}{*}{TG} & 3 & 2009 & \href{https://raumentwicklung.tg.ch/themen/kantonaler-richtplan.html/4211}{Online} & DE \\
 & 2 &  & † & DE \\
 & 1 &  & ‡ & DE \\ \hline

\multirow{3}{*}{UR} & 3 & 2012 & \href{https://www.ur.ch/dienstleistungen/3687}{Online} & DE \\
 & 2 &  & † & DE \\
 & 1 &  & ‡ & DE \\ \hline

\multirow{3}{*}{VD} & 3 & 2008 & \href{https://www.vd.ch/territoire-et-construction/amenagement-du-territoire/plan-directeur-cantonal/}{Online} & FR \\
 & 2 &  & † & FR \\
 & 1 &  & ‡ & FR \\ \hline

\multirow{3}{*}{VS} & 3 & 2019 & \href{https://www.vs.ch/fr/web/sdt/plan-directeur-cantonal-2019}{Online} & FR \\
 & 2 &  & † & FR \\
 & 1 &  & ‡ & FR \\ \hline

\multirow{3}{*}{ZG} & 3 & 2004 & \href{https://zg.ch/de/planen-bauen/raumplanung/richtplanung#archiv}{Online} & DE \\
 & 2 &  & † & DE \\
 & 1 &  & ‡ & DE \\ \hline

\multirow{3}{*}{ZH} & 3 & 2014 & \href{https://www.zh.ch/de/planen-bauen/raumplanung/richtplaene.html}{Online} & DE \\
 & 2 &  & † & DE \\
 & 1 &  & ‡ & DE \\ \hline

\caption{Cantonal master plans (part 2).}
\end{longtable}

\end{minipage}



\textbf{Legend:}
\begin{itemize}
    \item[†] Only available in paper format (physical consultation) and thus excluded from this study
    \item[‡] Sent upon request by the cantonal service.
\end{itemize}
}





The raw text was extracted from the pdf using a bash one liner

\begin{verbatim}
for f in *.pdf *.PDF; do
  [ -f "$f" ] && pdftotext "$f" "${f%.*}.txt"
done
\end{verbatim}


I then imported the .txt into TXM : https://txm.gitpages.huma-num.fr/textometrie/

TXM use the treetagger module to tag each word in the corpus. (https://www.cis.uni-muenchen.de/~schmid/tools/TreeTagger/) \cite{serge_txm_2010}

I used the French and German module for the French and German dataset.

I used TXM to understand a bit the corpus but in an attempt to make this more reproductible, I translated the useful query into python and made a script to renegerate my result.

I look at hit against a custom keyword list with word mapped to motilité and temps perçu

I did an AFC on the whole corpus to see if there are differences between documents but it did not yield any useful result. the figures are in annexes

I also did an AFC with the hit against time and moitilité